% \iffalse meta-comment
% Copyright (C) 2026 Russell Wright Helder
% \fi
%
% \iffalse
%<*driver>
\ProvidesFile{rhelder3.dtx}[2026/07/18 v0.1.0 Personal document class]
\documentclass{ltxdoc}

\usepackage{biblatex}
\usepackage{xfrac}

\addbibresource{rhelder3.bib}

\NewDocElement[
  envlike,
  idxtype=opt.,
  printtype=\textit{opt.}
]{Opt}{opt}

\NewDocElement[
  macrolike,
  idxtype=dimen,
  printtype=\textit{ldimen}
]{LDimen}{ldimen}

\NewDocElement[
  macrolike,
  idxtype=skip,
  printtype=\textit{lskip}
]{LSkip}{lskip}

\NewDocElement[
  envlike,
  idxtype=counter,
  printtype=\textit{counter}
]{LCounter}{lcounter}

\NewDocumentCommand{\acro}{m}{\textsc{\MakeLowercase{#1}}}
\NewDocumentCommand{\bool}{m}{\texttt{#1}}
\NewDocumentCommand{\env}{m}{\texttt{#1}}
\NewDocumentCommand{\hook}{m}{\texttt{#1}}
\NewDocumentCommand{\num}{m}{#1}

\begin{document}
  \DocInput{rhelder3.dtx}
\end{document}
%</driver>
% \fi
%
% \GetFileInfo{rhelder3.dtx}
%
% \title{The \cls{rhelder3} Class}
% \author{Russell Wright Helder}
% \date{\filedate \quad \fileversion}
%
% \maketitle
%
% \section{Identification}
%
% Define a prefix for internal commands (\pkg{docstrip} will replace
% \texttt{@@} with this prefix when generating the class file).
%    \begin{macrocode}
%<@@=rh>
%    \end{macrocode}
% Check that we are running on top of \LaTeXe{}
% and are using a recent version of the kernel
% that contains the L3 programming layer,
% \pkg{xparse},
% and the command \cs{ProcessKeyOptions}.
%    \begin{macrocode}
\NeedsTeXFormat{LaTeX2e}[2022-06-01]
%    \end{macrocode}
% Identify the class file and enable \pkg{expl3} syntax.
%    \begin{macrocode}
\ProvidesExplClass{rhelder3}{2026/07/18}{0.1.0}{Personal document class}
%    \end{macrocode}

% \section{Options}
%
% \begin{opt}{showgrid}
%     Boolean (default \bool{false}).
%     Show baseline grid. Useful for document design.
%    \begin{macrocode}
\keys_define:nn {rh}{
  showgrid .code:n =
    \RequirePackage{tikz}
    \dim_new:N \l_@@_to_first_baseline
    \dim_new:N \l_@@_yshift

    \hook_gput_code:nnn {begindocument}{.}{\@@_show_baseline_grid:},
}

\ProcessKeyOptions[rh]
%    \end{macrocode}
% \end{opt}

% \section{Horizontal Motion}
%
% \subsection{Measure}
%
% \begin{macro}{\SetTextwidth}
%    \marg{integer}\\
%    Sets \cs{textwidth}
%    so that the line will contain
%    an average of roughly \meta{integer} characters.
%    This formula for calculating the line width
%    is extracted from the copyfitting table at
%    \autocite[28]{bringhurst1992ElementsTypographic}.
%    \begin{macrocode}
\NewDocumentCommand{\SetTextwidth}{m}{\@@_set_textwidth:n {#1}}

\dim_new:N \l_@@_fontwidth

\cs_new:Npn \@@_set_textwidth:n #1 {
  \settowidth{\l_@@_fontwidth}{abcdefghijklmnopqrstuvwxyz}
  \geometry{
    textwidth=\dim_eval:n {
      (0.0325\l_@@_fontwidth + 0.378pt) * #1
    }
  }
}
%    \end{macrocode}
% \end{macro}

% \section{Leading}
%
% \begin{macro}{\normalsize}
%    Sets the font size and leading.
%    Also set \cs{parskip} equal to one lead,
%    and set generic list parameters.
%    \begin{macrocode}
\RenewDocumentCommand{\normalsize}{}{
  \@setfontsize\normalsize{12pt}{15pt}
  \setlength\parindent{\baselineskip}
  \let\@listi\@listI
}
%    \end{macrocode}
% \end{macro}

% \begin{macro}{\large}
%    Set other font sizes.
%    \begin{macrocode}
\NewDocumentCommand{\large}{}{\@setfontsize\large{14pt}{18pt}}
%    \end{macrocode}
% \end{macro}

% Initialize font size.
%    \begin{macrocode}
\normalsize
%    \end{macrocode}
%
% \begin{lskip}{\topskip}
%    Set leading for the first line on the page.
%    \begin{macrocode}
\setlength\topskip{12pt}
%    \end{macrocode}
% \end{lskip}

% \section{Paragraphs}
%
% \begin{ldimen}{\parindent}
%    Initialize \cs{parindent} to \cs{baselineskip},
%    although this is not strictly necessary
%    (\cs{parindent} will be set to \cs{baselineskip}
%    when \cs{normalsize} is called).
%    \begin{macrocode}
\setlength\parindent{\baselineskip}
%    \end{macrocode}
% \end{ldimen}
%
% \begin{lskip}{\parskip}
%    Don't put any stretchable space between paragraphs,
%    so that the text snaps to baseline.
%    \begin{macrocode}
\setlength\parskip{0pt}
%    \end{macrocode}
% \end{lskip}

% \subsection{Baseline Grid}
%
% Define the macro to be expanded if \opt{showgrid} is set.
% Useful for document design.
%    \begin{macrocode}
\cs_new:Npn \@@_show_baseline_grid: {
  \dim_set:Nn \l_@@_to_first_baseline {\Gm@tmargin + \topskip}
  \dim_set:Nn \l_@@_yshift {
    \fp_eval:n {
      \fp_to_int:n {\l_@@_to_first_baseline / \baselineskip} -
      \l_@@_to_first_baseline / \baselineskip
    }\baselineskip
  }

  \hook_gput_code:nnn {shipout/background}{.}{
    \put (0, 0){
      \begin{tikzpicture}[remember~picture, overlay]
        \draw [ystep=\baselineskip, xstep=0, help~lines]
          (0, 0) [yshift=\l_@@_yshift] grid +(\pagewidth, -\pageheight);
      \end{tikzpicture}
    }
  }
}
%    \end{macrocode}

% \section{Blocks}
%
% \subsection{Parameters of Generic \env{list} Environment}
%
% \begin{ldimen}{\leftmargin, \leftmargini}
%    Align lists with the surrounding paragraph indent.
%    \begin{macrocode}
\setlength\leftmargin{\parindent}
\setlength\leftmargini{\parindent}
%    \end{macrocode}
% \end{ldimen}

% \begin{macro}{\@listi, \@listI}
%    Don't put any stretchable space between list elements,
%    and set \cs{topsep} to half a \cs{baselineskip},
%    so that the text will snap back to baseline after the end of the list.
%    \begin{macrocode}
\def\@listi{
  \setlength\leftmargin{\leftmargini}
  \setlength\topsep{0.5\baselineskip}
  \setlength\parsep{0pt}
  \setlength\itemsep{0pt}
}
\let\@listI\@listi

\@listi
%    \end{macrocode}
% \end{macro}

% \subsection{Block Environments Based on \env{list}}
%
% \begin{env}{quotation, quote}
%    Define environments for setting block quotations.
%    Define both \env{quotation} and \env{quote},
%    as in the standard classes,
%    in case any packages expect one or the other.
%    \begin{macrocode}
\NewDocumentEnvironment{quotation}{}{
  \list{}{
    \setlength\itemindent{0pt}
    \setlength\listparindent{\parindent}
    \setlength\rightmargin{\leftmargin}
  }
  \item\relax
}{
  \endlist
}

\NewEnvironmentCopy{quote}{quotation}
%    \end{macrocode}
% \end{env}

% \section{Headings \& Subheads}
%
% \begin{macro}{
%   \section,
%   \subsection,
%   \subsubsection,
%   \paragraph,
%   \subparagraph
% }
%    Define all of the headings and subheads defined in the standard classes.
%    \begin{macrocode}
\newcommand\section{
  \@startsection{section}{1}
  {0pt}
  {-\fp_to_dim:n {(45 - 18) / 2.701 * 1.701}}
  {\fp_to_dim:n {(45 - 18) / 2.701 * 1}}
  {
    \savepos\RecordProperties{rh@ypos@\thesection}{ypos}
    \headfamily\scshape\large
  }
}

\newcommand\subsection{
  \@startsection{subsection}{2}
  {0pt}
  {-\fp_to_dim:n {(30 - 15) / 2.701 * 1.701}}
  {\fp_to_dim:n {(30 - 15) / 2.701 * 1}}
  {
    \savepos\RecordProperties{rh@ypos@\thesubsection}{ypos}
    \rmfamily\scshape\normalsize
  }
}

\newcommand\subsubsection{
  \@startsection{subsubsection}{3}
  {0pt}
  {-\fp_to_dim:n {(30 - 15) / 2.701 * 1.701}}
  {\fp_to_dim:n {(30 - 15) / 2.701 * 1}}
  {
    \savepos\RecordProperties{rh@ypos@\thesubsubsection}{ypos}
    \rmfamily\itshape\normalsize
  }
}

\newcommand\paragraph{
  \@startsection{paragraph}{4}
  {0pt}
  {\baselineskip}
  {-1em}
  {\rmfamily\bfseries\normalsize}
}

\newcommand\subparagraph{
  \@startsection{subparagraph}{5}
  {0pt}
  {\baselineskip}
  {-1em}
  {\rmfamily\itshape\normalsize}
}
%    \end{macrocode}
% \end{macro}

% \begin{lcounter}{secnumdepth}
%    Number all headings down to \cs{subsubsection}.
%    \begin{macrocode}
\setcounter{secnumdepth}{3}
%    \end{macrocode}
% \end{lcounter}

% \begin{lcounter}{
%   section,
%   subsection,
%   subsubsection,
%   paragraph,
%   subparagraph,
% }
%    Define counters used for section numbers.
%    \begin{macrocode}
\newcounter{section}
\newcounter{subsection}[section]
\newcounter{subsubsection}[subsection]
\newcounter{paragraph}[subsubsection]
\newcounter{subparagraph}[paragraph]
%    \end{macrocode}
% \end{lcounter}

% \begin{macro}{
%   \thesection,
%   \thesubsection,
%   \thesubsubsection,
%   \theparagraph,
%   \thesubparagraph
% }
%    Format section numbers.
%    \begin{macrocode}
\renewcommand\thesection{\arabic{section}}
\renewcommand\thesubsection{\thesection.\arabic{subsection}}
\renewcommand\thesubsubsection{\thesubsection.\arabic{subsubsection}}
\renewcommand\theparagraph{\thesubsubsection.\arabic{paragraph}}
\renewcommand\thesubparagraph{\theparagraph.\arabic{subparagraph}}

\def\@seccntformat#1{\normalfont\csname the#1\endcsname\quad}
%    \end{macrocode}
% \end{macro}

% If a display heading or subhead is on the first line of the page,
% shift the heading or subhead down from the first baseline
% so that the subsequent text snaps to the baseline grid.
%    \begin{macrocode}
\hook_gput_code:nnn {cmd/section/before}{.}{
  % The section counter hasn't been incremented yet. Temporarily increment in
  % order to reference the position of the right section.
  \addtocounter{section}{1}
  \@@_shift_if_at_top:n {section}{18pt}
  \addtocounter{section}{-1}
}

\hook_gput_code:nnn {cmd/subsection/before}{.}{
  % The subsection counter hasn't been incremented yet. Temporarily increment
  % in order to reference the position of the right section.
  \addtocounter{subsection}{1}
  \@@_shift_if_at_top:n {subsection}{18pt}
  \addtocounter{subsection}{-1}
}

\hook_gput_code:nnn {cmd/subsubsection/before}{.}{
  % The subsubsection counter hasn't been incremented yet. Temporarily
  % increment in order to reference the position of the right section.
  \addtocounter{subsubsection}{1}
  \@@_shift_if_at_top:n {subsubsection}{18pt}
  \addtocounter{subsubsection}{-1}
}
%    \end{macrocode}

% Get the y-coordinate of the top of the text block.
%    \begin{macrocode}
\hook_gput_code:nnn {begindocument/end}{.}{
  \savepos\RecordProperties{rh@ypos@topofpage}{ypos}
}
%    \end{macrocode}

% Define the macro that is expanded in the \hook{cmd/section/before} hook.
% It accepts the \meta{csname} of the heading
% plus the \meta{skip} by which to shift the header
% as arguments.
%
% The macro works by inserting a dimensionless rule
% (before the heading,
% since the macro is expanded in the \hook{cmd/section/before} hook),
% in order to make it possible to insert glue
% between the heading and the top of the page.
% As a consequence of this,
% \cs{topskip} is inserted;
% neutralize this effect by inserting a negative vertical skip
% equal to \cs{topskip}.
% Finally,
% add an additional vertical skip
% specified by the argument \meta{skip}.
%    \begin{macrocode}
\int_new:N \l_@@_ypos_top
\int_new:N \l_@@_ypos_heading

\cs_new:Npn \@@_shift_if_at_top:n #1#2 {
  \int_set:Nn \l_@@_ypos_top {
    \RefProperty{rh@ypos@topofpage}{ypos}
  }
  \int_set:Nn \l_@@_ypos_heading {
    \RefProperty{rh@ypos@\cs:w the #1\cs_end:}{ypos}
  }

  \bool_if:nT {
    \int_compare_p:nNn {\l_@@_ypos_heading} = {\l_@@_ypos_top}
    || \int_compare_p:nNn {\l_@@_ypos_heading} > {\l_@@_ypos_top}
  }{
    \rule{0pt}{0pt}
    \vspace{-\topskip}
    \vspace{-#2}
  }
}
%    \end{macrocode}

% \section{Load Packages}
%
%    \begin{macrocode}
\RequirePackage{fontspec}
%    \end{macrocode}
%
% Abstract from \LaTeX{}'s counterintuitive layout parameters
% with the \pkg{geometry} package.
%    \begin{macrocode}
\RequirePackage[letterpaper]{geometry}
%    \end{macrocode}

% \section{Font}
%
% Set \acro{EB} Garamond as the default Roman font.
%    \begin{macrocode}
\setmainfont[
  Extension = .otf,
  UprightFont = *-Regular,
  ItalicFont = *-Italic,
  BoldFont = *-SemiBold,
  BoldItalicFont = *-SemiBoldItalic,
%    \end{macrocode}
%
% Use text figures by default.
%    \begin{macrocode}
  Numbers = {Proportional, OldStyle},
%    \end{macrocode}
%
% Set up Swash font.
%    \begin{macrocode}
  FontFace = {m}{sw}{Font=*-Italic, Style=Swash},
  FontFace = {b}{sw}{Font=*-SemiBoldItalic, Style=Swash},
  FontFace = {m}{scsw}{Font=*-Italic, Letters=SmallCaps, Style=Swash},
%    \end{macrocode}
%
% Letterspace small caps.
%    \begin{macrocode}
  SmallCapsFeatures = {LetterSpace=4.0},
]{EBGaramond}
%    \end{macrocode}

% \begin{macro}{\headfamily}
%    Use Cormorant Garamond,
%    which is lighter than \acro{EB} Garamond,
%    as a display font.
%    \begin{macrocode}
\newfontfamily{\headfamily}{CormorantGaramond}
%    \end{macrocode}
% \end{macro}

% \begin{macro}{\initfamily, \textinit}
%    Switches to \acro{EB} Garamond's font face
%    with decorative initials.
%    \begin{macrocode}
\newfontface\initfamily{EBGaramond-Initials}
\DeclareTextFontCommand{\textinit}{\initfamily}
%    \end{macrocode}
% \end{macro}

% \begin{macro}{
%   \textfigures,
%   \liningfigures,
%   \tabularfigures,
%   \proportionalfigures
% }
%    Switch between text figures and lining figures,
%    tabular figures and proportional figures.
%    \begin{macrocode}
\DeclareTextFontCommand{\textfigures}{
  \addfontfeatures{Numbers=OldStyle}
}
\DeclareTextFontCommand{\liningfigures}{
  \addfontfeatures{Numbers=Lining}
}
\DeclareTextFontCommand{\tabularfigures}{
  \addfontfeatures{Numbers=Tabular}
}
\DeclareTextFontCommand{\proportionalfigures}{
  \addfontfeatures{Numbers=Proportional}
}
%    \end{macrocode}
% \end{macro}

% \begin{macro}{\sufigures, \textsu}
%    Typeset superior glyphs.
%    \begin{macrocode}
\newcommand\sufigures{\addfontfeatures{VerticalPosition=Superior}}
\DeclareTextFontCommand{\textsu}{\sufigures}
%    \end{macrocode}
% \end{macro}

% \begin{macro}{\infigures, \textin}
%    Typeset inferior glyphs.
%    \begin{macrocode}
\newcommand\infigures{\addfontfeatures{VerticalPosition=ScientificInferior}}
\DeclareTextFontCommand{\textin}{\infigures}
%    \end{macrocode}
% \end{macro}

% \begin{macro}{\textsw}
%    \marg{text}\\
%    Typeset \meta{text} in Swash shape.
%    \begin{macrocode}
\DeclareTextFontCommand{\textsw}{\swshape}
%    \end{macrocode}
% \end{macro}

% \begin{macro}{\MakeUppercase}
%    Letterspace capitals.
%    \begin{macrocode}
\NewCommandCopy\MakeUppercaseCopy\MakeUppercase
\RenewDocumentCommand{\MakeUppercase}{m}{
  \begingroup
    \addfontfeature{LetterSpace=9.0}\MakeUppercaseCopy{#1}
  \endgroup
}
%    \end{macrocode}
% \end{macro}

% \section{The Page}
%
% \subsection{Text Block}
%
% \begin{macro}{\SetTextheight}
%    \marg{ratio}\\
%    Set \(\sfrac{\cs{textheight}}{\cs{textwidth}}\)
%    to be equal to \meta{ratio}.
%    \begin{macrocode}
\NewDocumentCommand{\SetTextheight}{m}{\@@_set_textheight:n {#1}}

\int_new:N \l_@@_num_lines
\dim_new:N \l_@@_tmp

\cs_new:Npn \@@_set_textheight:n #1 {
  \int_set:Nn \l_@@_num_lines {
    \fp_to_int:n {
      (#1\textwidth - \topskip) / \baselineskip
    }
  }
  \geometry{
    textheight=\dim_eval:n {
      \topskip + \l_@@_num_lines\baselineskip
    }
  }
}
%    \end{macrocode}
% \end{macro}

% Choose a proportion between text height and width
% that appealingly contrasts with the letter page height and width.
% A ratio of \num{1.701}
% (tall pentagon; see \autocite[150]{bringhurst1992ElementsTypographic})
% seems to work well.
%    \begin{macrocode}
\SetTextwidth{70}
\SetTextheight{1.701}
%    \end{macrocode}
%
% Set the bottom margin equal to the horizontal margin.
% This creates an appealing and undramatic
% approximately \num{2:3} ratio
% between the top and bottom margins,
% at least when the line accommodates 70 characters on average.
%    \begin{macrocode}
\geometry{
  hcentering,
  bmargin=\dim_eval:n {(\paperwidth - \textwidth) / 2}
}
%    \end{macrocode}

% \subsection{Page Numbering}
%
%    \begin{macrocode}
\geometry{footskip=3\baselineskip}
\pagenumbering{arabic}
\pagestyle{plain}
%    \end{macrocode}
%
% \printbibliography
